% SKY2100 -- Exercise 7 (self-contained article; no external .sty needed)
\documentclass[11pt]{article}

% ---- Packages ------------------------------------------------------
\usepackage[a4paper,margin=2.2cm]{geometry}
\usepackage[T1]{fontenc}
\usepackage[utf8]{inputenc}
\usepackage{xcolor}
\usepackage{enumitem}
\usepackage{pgffor}
\usepackage{amssymb}
\usepackage{array}

% ---- Colours -------------------------------------------------------
\definecolor{kred}{HTML}{D81E2E}
\definecolor{kgrey}{HTML}{595959}

\usepackage[colorlinks=true,urlcolor=kred,linkcolor=kred]{hyperref}

% ---- Worksheet macros (inlined) ------------------------------------
\newcommand{\wbsection}[1]{\par\vspace{11pt}\noindent
  {\large\bfseries\color{kred}#1}\par\vspace{2pt}
  {\color{kred!35}\hrule height 1pt}\par\vspace{6pt}}

\newcommand{\task}[1]{\par\vspace{4pt}\noindent
  \textcolor{kred}{\textbf{$\triangleright$}}~#1\par\vspace{2pt}}

% Answers are discussed verbally -- no writing rules, just a small gap.
\newcommand{\answerlines}[1]{\par\vspace{0.4em}}

\newcommand{\boxlabel}[2]{\par\vspace{3pt}\noindent
  \fbox{\parbox{\dimexpr\linewidth-2\fboxsep-2\fboxrule}{\textbf{#1:}\enspace #2}}\par\vspace{3pt}}
\newcommand{\evidence}[1]{\boxlabel{Evidence to capture}{#1}}
\newcommand{\note}[1]{\boxlabel{Note}{#1}}
\newcommand{\caution}[1]{\boxlabel{Caution}{#1}}

% ---- Aha box + no italics (added) ----------------------------------
\newcommand{\aha}[1]{\par\vspace{5pt}\noindent
  \setlength{\fboxrule}{1.2pt}%
  {\color{kred}\fbox{\normalcolor\parbox{\dimexpr\linewidth-2\fboxsep-2\fboxrule}{%
     {\bfseries\color{kred}AHA.}\enspace #1}}}%
  \setlength{\fboxrule}{0.4pt}\par\vspace{5pt}}
\renewcommand{\emph}[1]{\textup{#1}}

\newcommand{\cmd}[1]{\texttt{#1}}
\newcommand{\cbox}{$\square$\enspace}

% ---- Hint and solutions divider -- same definitions as in kristiania-workbook.sty ----
\newcommand{\hint}[1]{\par\vspace{1pt}{\footnotesize\color{kgrey}\textbf{Hint:}\enspace #1}\par\vspace{2pt}}
\newcommand{\solheader}{\par\vspace{9pt}\noindent{\color{kred}\rule{\linewidth}{1.2pt}}\par\vspace{4pt}}

\newcommand{\signoff}{\par\vspace{9pt}\noindent
  {\color{kred}\rule{\linewidth}{0.4pt}}\par\vspace{3pt}
  \noindent{\footnotesize\color{kgrey}\textbf{Progress check:} Discuss your answers in the group, then
  talk the teacher through them verbally at the next session only to track progress; nothing is handed
  in, signed or graded.}\par}

\setlist[enumerate]{itemsep=3pt,topsep=2pt,parsep=0pt}
\setlist[itemize]{itemsep=2pt,topsep=2pt,parsep=0pt}
\setlength{\parindent}{0pt}
\setlength{\parskip}{2pt}

\begin{document}
\begin{center}
  {\LARGE\bfseries Harden your service \& control access}\\[3pt]
  {\large Session 7 -- Hardening \& access control}\\[5pt]
  {\small Exercise 7 \textperiodcentered\ Session 7 lab \textperiodcentered\ Estimated time: 90--120 min}
\end{center}
\vspace{2pt}\hrule\vspace{1.1em}

Write your answers in the answer document \cmd{SKY2100\_Exercise7\_Answers.docx}. Last week you put Nextcloud behind HTTPS. The service is
encrypted but still exposed: default settings, an open SSH port, a single admin account. Today you
\emph{shrink the attack surface} and decide \emph{who} may get in. You will reduce running services,
lock down the firewall and SSH, add brute-force protection, and apply least privilege inside the app.

\noindent\textbf{Caution:} Work only on your \emph{own} VM. Before you disable SSH password login,
make sure key-based login already works in a second terminal -- so you cannot lock yourself out.

\wbsection{1\quad Map the attack surface}
\task{List everything that is listening on your VM and who owns each port:}
\begin{quote}\ttfamily
ss -tulpn
\end{quote}
\task{Which ports are exposed, and which do you actually need for an HTTPS Nextcloud you administer
over SSH? Mark any you can switch off.}

\wbsection{2\quad Default-deny firewall (UFW)}
\task{Configure UFW to deny everything inbound, then allow only what you need:}
\begin{quote}\ttfamily
sudo ufw --force reset\hfill\# start clean: the S4 rules (22 and 80 open to all) would otherwise stay\\
sudo ufw default deny incoming\\
sudo ufw default allow outgoing\\
sudo ufw allow 443/tcp\\
sudo ufw allow from <your-ip> to any port 22\\
sudo ufw enable \&\& sudo ufw status verbose
\end{quote}
\task{Why is ``default deny, then allow'' safer than ``allow the bad ports''? (one or two lines)}

\wbsection{3\quad Harden SSH}
\task{Set up key-based login (if not already), then in \cmd{/etc/ssh/sshd\_config} set:}
\begin{quote}\ttfamily
PasswordAuthentication no\\
PermitRootLogin no
\end{quote}
Reload with \cmd{sudo systemctl reload ssh}. \textbf{Keep your current session open} and confirm a new
key-based login works before closing it.
\task{Why disable password authentication and root login specifically? Link each to a threat.}

\wbsection{4\quad Brute-force protection (fail2ban)}
\task{Install and enable fail2ban for SSH:}
\begin{quote}\ttfamily
sudo apt install -y fail2ban\\
sudo systemctl enable --now fail2ban\\
sudo fail2ban-client status sshd
\end{quote}
\task{Trigger a few failed logins from another machine/IP, then check the jail again. What changed?}

\wbsection{5\quad Least privilege inside the app}
\task{In Nextcloud, create a \emph{non-admin} user and give it only the access it needs. Do \emph{not}
use the admin account for everyday work.}
\task{Enable two-factor authentication (2FA) on the admin account.}
\task{State one thing the admin can do that your new user cannot -- and why that is least privilege.}

\wbsection{6\quad Cloud transfer: Microsoft Learn}
\task{Work through ``Secure your Azure resources with Azure role-based access control (Azure RBAC)''.}
\par\noindent\mbox{\href{https://learn.microsoft.com/training/modules/secure-azure-resources-with-rbac/}{learn.microsoft.com/training/modules/secure-azure-resources-with-rbac}}
\task{In Azure RBAC you assign a \textbf{role} to an \textbf{identity} at a \textbf{scope}. Name the
four scope levels from broadest to narrowest, and say which you would choose for least privilege.}

\wbsection{7\quad Azure reflection}
\task{Your UFW rules protect one VM. Which Azure construct enforces the \emph{same} default-deny idea
across your cloud resources, and at which layer (host vs.\ network)?}
\task{CSA calls IAM ``the new perimeter''. In your own words, why does identity matter \emph{more} in
the cloud than the network perimeter did on-premises?}

\aha{Last week the service was encrypted, yet today you still had work to do: an open SSH port, a single
all-powerful admin, default settings. Encryption did not make the service safe, because reachability and
access are separate problems from confidentiality on the wire. Notice where the real boundary sits now.
On-premises you could draw a line around the building and trust what was inside, but your cloud service
has no inside; anyone who reaches it is judged by who they are, not where they connect from. That is why
identity is the new perimeter, and why least privilege, MFA and a locked-down SSH matter more than any
network wall.}

\wbsection{8\quad Knowledge check}
\begin{enumerate}
\item Define \textbf{attack surface} and give two ways hardening reduces it.
\hint{Every point an attacker could enter or extract data — hardening shrinks it.}
\item Why can a cloud tenant \emph{not} rely on traditional perimeter security? (Comer)
\hint{You share physical servers and links with other tenants — no hard perimeter.}
\item State the difference between \textbf{authentication} and \textbf{authorisation}.
\hint{One proves who you are; one decides what you may do.}
\item Name the three authentication factors and say what \textbf{MFA} combines.
\hint{Something you know, have, are — and combining two different ones.}
\item Distinguish \textbf{RBAC} from \textbf{ABAC} with one example each.
\hint{Permissions by role, versus permissions by attributes and conditions.}
\item State the \textbf{principle of least privilege} in one sentence.
\hint{The minimum access an identity needs, and nothing more.}
\item CSA defines an \textbf{entitlement} as a rule over which four things?
\hint{A rule: which entity, which actions, which resources, under which conditions.}
\item What does \textbf{zero trust} mean for how each request is handled?
\hint{Trust nothing by default; verify every request instead of a one-time gate.}
\item What is \textbf{single sign-on (SSO)}, and what problem does it solve?
\hint{One login across services from a central identity system — no password sprawl.}
\item What is \textbf{PAM}, and why are privileged accounts treated specially?
\hint{Identity for admin accounts, which can do the most damage if taken.}
\item In Azure RBAC, what three things make up a role assignment?
\hint{Identity, role, and the scope it applies to.}
\item Name two \textbf{common cloud access-control mistakes} and the fix for each.
\hint{Admin roles without MFA, long-lived static keys, over-broad roles.}
\end{enumerate}

\signoff

\clearpage
\solheader
{\LARGE\bfseries\color{kred} Solutions}\\[2pt]
{\small Work through the lab first, then check here. Model answers are indicative — equivalent reasoning is fine.}\par\vspace{8pt}
\noindent\textbf{Note:} Model answers are indicative -- accept equivalent reasoning. Multiple-choice
answers are in \textbf{bold}.

\wbsection{Knowledge check}
\begin{enumerate}
\item The \textbf{attack surface} is the sum of all points where an attacker could try to enter or
      extract data -- open ports, services, accounts, APIs, dependencies. Hardening reduces it by, e.g.,
      removing unused services/packages and closing ports (default-deny firewall).
\item In the cloud, VMs and containers share physical servers and links with other tenants, so attacks
      can arise \emph{from within} the facility; there is no clear inside/outside, so traditional
      perimeter security cannot define the network boundary (Comer p.~236).
\item \textbf{Authentication} verifies \emph{who} you are (identity of a user/process/device).
      \textbf{Authorisation} decides \emph{what} you may do (access control via granted permissions).
      Authenticate first, then authorise (CSA p.~105).
\item Factors: \textbf{something you know} (password), \textbf{something you have} (phone/token),
      \textbf{something you are} (biometric). \textbf{MFA} combines two or more, so a stolen password
      alone is not enough.
\item \textbf{RBAC} attaches permissions to roles and assigns roles to users (e.g.\ ``editor''
      $\rightarrow$ edit files). \textbf{ABAC} grants access based on attributes/conditions (e.g.\
      only with MFA on, only from a managed device, only if a resource tag matches) (CSA p.~105).
\item \textbf{Least privilege}: give every identity the minimum access it needs to do its job, and no
      more (Comer p.~236--237; CSA p.~119).
\item An \textbf{entitlement} is a rule: this \textbf{entity} may perform \textbf{these actions} on
      \textbf{these resources} under \textbf{these conditions} (with these attributes) (CSA p.~119).
\item \textbf{Zero trust}: trust nothing by default; validate \emph{every} request using identity (and
      attributes/risk), instead of granting blanket access after a single login (Comer p.~237; CSA
      p.~120).
\item \textbf{Single sign-on (SSO)}: one set of credentials for all services via a central identity
      system; it solves the problem of separate logins per service and lets a user authenticate once
      per task (Comer p.~238).
\item \textbf{PAM (Privileged Access Management)} handles identity for administrative/superuser
      accounts. They are treated specially because forging one lets an attacker cause serious damage;
      PAM limits each admin to the systems they manage and logs all (and failed) accesses (Comer
      p.~238--239).
\item A role assignment = an \textbf{identity} (user/group/service principal/managed identity) +
      a \textbf{role} (e.g.\ Reader, Contributor) + a \textbf{scope} (management group / subscription /
      resource group / resource).
\item Any two of: admin roles without MFA; long-lived static access keys; over-broad roles (Owner where
      Reader suffices); public-by-accident storage; management plane exposed with weak auth. Fix:
      enforce MFA; rotate/eliminate static secrets (use managed identities); apply least-privilege
      roles; close public access; restrict and protect the management plane (CSA p.~104, 120).
\end{enumerate}

\wbsection{Lab checkpoints}
\begin{itemize}
\item \textbf{Surface:} \cmd{ss -tulpn} should show only what is needed (443 for HTTPS, 22 for SSH).
      Anything else -- old test servers, databases bound to \cmd{0.0.0.0} -- is a candidate to disable
      or bind to localhost.
\item \textbf{Firewall:} \cmd{ufw status verbose} must show \emph{default deny incoming} plus explicit
      allows. ``Default deny then allow'' fails safe -- an unforeseen service is blocked unless you
      opened it; ``block the bad ports'' fails open.
\item \textbf{SSH:} key-based login works; \cmd{PasswordAuthentication no} defeats password
      guessing/brute force; \cmd{PermitRootLogin no} removes the highest-value target (use a normal
      user + \cmd{sudo}). Verify a new session \emph{before} closing the old one.
\item \textbf{fail2ban:} after a few failed attempts, \cmd{fail2ban-client status sshd} shows failed/
      banned counts -- continuous brute force becomes a few blocked tries. This is detection feeding
      response (Comer p.~238: record failed logins to track attackers).
\item \textbf{App least privilege:} a non-admin user cannot change server settings, manage users, or
      install apps; admin has 2FA enabled. Everyday work should not use the admin account.
\end{itemize}

\wbsection{Cloud transfer \& reflection}
\begin{itemize}
\item \textbf{RBAC scope (broadest$\rightarrow$narrowest):} management group, subscription, resource
      group, resource. For least privilege, assign at the \emph{narrowest} scope that still lets the
      task succeed.
\item \textbf{Azure default-deny:} a \textbf{Network Security Group (NSG)} enforces allow/deny by port/
      protocol/source at the \emph{network} layer around resources (the host UFW works at the host
      layer). Both should be default-deny.
\item \textbf{IAM as the new perimeter:} cloud consolidates data-centre controls into Internet-facing
      consoles and APIs protected mainly by identity; with no network wall, identity is what stands
      between an attacker and everything. CSA notes the vast majority of cloud-native breaches start
      with IAM failures (CSA p.~103--104).
\end{itemize}

\wbsection{References}
CSA Security Guidance v5, \textbf{Domain~5: Identity \& Access Management} (p.~103--120, verified) and
\textbf{Domain~7: Infrastructure \& Networking} (p.~147--148, verified). Comer, \emph{The Cloud
Computing Book}, Ch.~17 -- Cloud Security \& Privacy: attack surface (p.~233--235), perimeter \& zero
trust (p.~236--237), Identity Management \& SSO (p.~238), PAM (p.~238--239), verified. MS Learn:
\emph{Secure your Azure resources with Azure RBAC} (verified, June 2026) -- role + identity + scope,
with an activity log of RBAC changes.

\end{document}
